\section*{Capítulo \ref{ch:b3:galois} --- Extensões de corpos e teoria de
Galois}

\begin{solution}{exo:b3:galois:1}
$\sqrt3 \notin \Q(\sqrt2)$: de $\sqrt3 = a + b\sqrt2$ ($a, b
\in \Q$), elevar ao quadrado dá $3 = a^2 + 2b^2 + 2ab\sqrt2$, logo $ab =
0$; $b = 0$
torna $\sqrt3$ racional, $a = 0$ dá $\sqrt6 = 2b
\in \Q$ --- ambos falsos
(argumentos usuais de fatoração em primos). Logo $[\Q(\sqrt2,\sqrt3) : \Q(\sqrt2)] = 2$ e a \omterm{thm:b3:galois:tower}{lei das torres}
dá \omterm{def:b3:galois:extension}{grau} $4$.

Seja $\gamma = \sqrt2 + \sqrt3$. Então $\gamma^2 = 5 + 2\sqrt6$ e
$(\gamma^2 - 5)^2 = 24$: $\gamma$ anula $X^4 - 10X^2 + 1$.
Além disso $\gamma^3 = 11\sqrt2 + 9\sqrt3$, logo $\gamma^3 - 9\gamma
= 2\sqrt2$: $\sqrt2 \in \Q(\gamma)$, e então $\sqrt3 = \gamma -
\sqrt2 \in \Q(\gamma)$: $\Q(\gamma) = \Q(\sqrt2,\sqrt3)$, de
\omterm{def:b3:galois:extension}{grau} $4$. A quártica que o anula, tendo o \omterm{def:b3:galois:extension}{grau} do \omterm{def:b3:galois:algebraic}{polinômio
minimal}, \emph{é} o \omterm{def:b3:galois:algebraic}{polinômio minimal}: $X^4 -
10X^2 + 1$ (em particular, ele é \omterm{def:b3:rings:divisibility}{irredutível} sobre $\Q$).
\end{solution}

\begin{solution}{exo:b3:galois:2}
As três raízes de $X^3 - 2$ em $\C$ são $\alpha, j\alpha,
j^2\alpha$ com $j = \eu^{2\iu\pi/3}$. A aplicação $\alpha \mapsto
j\alpha$ define um $\Q$-mergulho $\Q(\alpha) \to \C$
(\cref{thm:b3:galois:simple}: ambos geram extensões de \omterm{def:b3:galois:extension}{grau} $3$
com o mesmo \omterm{def:b3:galois:algebraic}{polinômio minimal}), cuja imagem $\Q(j\alpha)
\not\subseteq \R$ difere de $\Q(\alpha) \subseteq \R$:
$\Q(\alpha)/\Q$ não é normal. O \omterm{thm:b3:galois:splitting}{corpo de decomposição} é $L =
\Q(\alpha, j)$, com $[L:\Q] = [L:\Q(\alpha)]\,[\Q(\alpha):\Q] =
2 \cdot 3 = 6$ ($j$ satisfaz $X^2 + X + 1$, \omterm{def:b3:rings:divisibility}{irredutível} sobre
o corpo real $\Q(\alpha)$). Um automorfismo de $\Q(\alpha)$
deve levar $\alpha$ a uma raiz de $X^3 - 2$ \emph{dentro de}
$\Q(\alpha) \subseteq \R$: só $\alpha$ serve, de modo que
$\operatorname{Aut}_\Q(\Q(\alpha)) = \{\mathrm{id}\}$, de ordem
$1 < 3$.
\end{solution}

\begin{solution}{exo:b3:galois:3}
$X^2 + 1$ não tem raiz em $\mathbb F_3$ ($0, 1, 2 \mapsto 1, 2,
2$), de modo que $\mathbb F_9 = \mathbb F_3[X]/(X^2+1)$ é um corpo com
$9$ elementos; escreva $\omega = \bar X$, $\omega^2 = -1$. O
grupo $\mathbb F_9^\times$ é cíclico de ordem $8$; $\omega$ tem
ordem $4$, mas $1 + \omega$ funciona: $(1+\omega)^2 = 1 + 2\omega +
\omega^2 = 2\omega$, $(1+\omega)^4 = 4\omega^2 = \omega^2 = -1
\ne 1$: ordem $8$.

Sobre $\mathbb F_2$ --- \omterm{def:b3:galois:extension}{grau} $1$: $X$, $X + 1$; \omterm{def:b3:galois:extension}{grau} $2$:
$X^2 + X + 1$ (as outras três quadráticas têm raízes); \omterm{def:b3:galois:extension}{grau}
$3$: $X^3 + X + 1$ e $X^3 + X^2 + 1$ (sem raízes em $\mathbb
F_2$; as outras seis cúbicas têm raízes). Verificação:
\[
(X^3{+}X{+}1)(X^3{+}X^2{+}1) = X^6 + X^5 + X^4 + X^3 + X^2 + X
+ 1,
\]
e $X(X{+}1)(X^6 + \dots + 1) = X(X^7 + 1) = X^8 + X = X^8 -
X$ sobre $\mathbb F_2$ --- exatamente os irredutíveis de \omterm{def:b3:galois:extension}{grau}
divisor de $3$, como o~\cref{exo:b3:galois:6} prevê (o \omterm{def:b3:galois:extension}{grau} $2$ está
ausente: $2 \nmid 3$).
\end{solution}

\begin{solution}{exo:b3:galois:4}
(a) Módulo $7$: as potências de $3$ são $3, 2, 6, 4, 5, 1$: ordem
$6$, um gerador; as raízes primitivas são os $3^k$ com
$\gcd(k, 6) = 1$: $3$ e $3^5 = 5$. Módulo $11$: potências de $2$:
$2, 4, 8, 5, 10, 9, 7, 3, 6, 1$: um gerador; raízes primitivas
$2^k$, $\gcd(k, 10) = 1$: $2, 2^3 = 8, 2^7 = 7, 2^9 = 6$.

(b) Escreva $x = g^k$ com $g$ um gerador
(\cref{thm:b3:galois:cyclic}). Então $x$ é um quadrado se, e somente se, $k$ é
par (os quadrados são os $g^{2l}$, e $g^{2l} = g^{k}$ se, e somente se, $k
\equiv 2l \bmod p-1$, \omterm{def:b3:groups:derived}{solúvel} se, e somente se, $k$ é par, sendo $p - 1$
par). E $x^{(p-1)/2} = g^{k(p-1)/2} = 1$ se, e somente se, $(p-1) \mid
k\frac{p-1}2$ se, e somente se, $k$ é par: as duas condições coincidem. Para $x =
-1 = g^{(p-1)/2}$: é um quadrado se, e somente se, $\frac{p-1}2$ é par, se, e somente se,
$p \equiv 1 \pmod 4$ --- o~\cref{pb:b3:rings:1} de novo.
\end{solution}

\begin{solution}{exo:b3:galois:5}
Dentro de $\bar{\mathbb F}_p$, $\mathbb F_{p^k} = \{x : x^{p^k} =
x\}$ é o conjunto dos pontos fixos de $F^k$. A interseção $\mathbb F_{p^m}
\cap \mathbb F_{p^n}$ é fixada por $F^m$ e $F^n$, logo por
$F^{\gcd(m,n)}$ ($\gcd = am + bn$: sobre um elemento fixo, $F^{am +
bn} = (F^m)^a(F^n)^b$ age trivialmente --- os expoentes podem ser tomados
positivos por periodicidade); assim ela está contida em $\mathbb
F_{p^{\gcd(m,n)}}$, que reciprocamente está contido em ambos
(\cref{thm:b3:galois:finitefields}(3)). O composto $\mathbb
F_{p^m}\mathbb F_{p^n}$: todo corpo que contém ambos tem \omterm{def:b3:galois:extension}{grau}
divisível por $m$ e $n$, logo por $\operatorname{lcm}(m,n)$; e
$\mathbb F_{p^{\operatorname{lcm}}}$ contém ambos: é o
composto. Para $m \mid n$: $\operatorname{Gal}(\mathbb
F_{p^n}/\mathbb F_{p^m})$ consiste nas potências de $F$ que fixam
$\mathbb F_{p^m}$, isto é, nas de $F^m$: cíclico de ordem $n/m$,
gerado por $F^m \colon x \mapsto x^{p^m}$ (ordem como em
\cref{thm:b3:galois:finitefields}(2)).
\end{solution}

\begin{solution}{exo:b3:galois:6}
$X^{q^n} - X$ é \omterm{def:b3:galois:separable}{separável} (derivada $-1$), com conjunto de raízes
$\mathbb F_{q^n}$. Seja $P$ mônico \omterm{def:b3:rings:divisibility}{irredutível} de \omterm{def:b3:galois:extension}{grau} $d$.
Se $d \mid n$: $\mathbb F_q[X]/(P) \cong \mathbb F_{q^d}
\subseteq \mathbb F_{q^n}$, de modo que $P$ tem uma raiz $\alpha \in
\mathbb F_{q^n}$; $\alpha^{q^n} = \alpha$, e $P =
\pi_\alpha$ divide $X^{q^n} - X$. Se $P \mid X^{q^n} - X$: uma
raiz $\alpha \in \mathbb F_{q^n}$ gera $\mathbb F_{q^d}
\subseteq \mathbb F_{q^n}$, logo $d \mid n$
(\cref{thm:b3:galois:finitefields}(3)). Irredutíveis distintos
são primos entre si e o produto é \omterm{def:b3:galois:separable}{separável}: cada $P$ aparece com
expoente exatamente $1$, e toda raiz de $X^{q^n}-X$ é raiz
de seu \omterm{def:b3:galois:algebraic}{polinômio minimal}: a fatoração vale. Comparando
\omterm{def:b3:galois:extension}{graus}: $q^n = \sum_{d\mid n} d\,I_d(q)$.

Consequentemente $I_1 = q$; $q^2 = I_1 + 2I_2$ dá $I_2 =
\frac{q^2 - q}2$; $q^3 = I_1 + 3I_3$ dá $I_3 = \frac{q^3 -
q}3$; $q^4 = I_1 + 2I_2 + 4I_4$ dá $I_4 = \frac{q^4 -
q^2}4$. Existência: $nI_n = q^n - \sum_{d \mid n,\, d < n}
dI_d \geq q^n - \sum_{d \leq n/2} q^d > q^n - q^{n/2 + 1} \geq
0$ para $n \geq 2$ (e $I_1 = q \geq 1$): $I_n \geq 1$ sempre.
\end{solution}

\begin{solution}{exo:b3:galois:7}
$L = \Q(\sqrt2, \sqrt3)$ é o \omterm{thm:b3:galois:splitting}{corpo de decomposição} de $(X^2 -
2)(X^2 - 3)$, \omterm{def:b3:galois:separable}{separável}: galoisiano de \omterm{def:b3:galois:extension}{grau} $4$
(\cref{exo:b3:galois:1}). Um automorfismo leva $\sqrt2 \mapsto
\pm\sqrt2$ e $\sqrt3 \mapsto \pm\sqrt3$: no máximo $4$ escolhas,
e $\abs G = 4$ realiza todas: $G \cong (\Z/2\Z)^2$, com
elementos $\mathrm{id}, \sigma (\sqrt2 \mapsto -\sqrt2), \tau
(\sqrt3\mapsto-\sqrt3), \sigma\tau$. Subgrupos de ordem $2$:
$\langle\sigma\rangle, \langle\tau\rangle,
\langle\sigma\tau\rangle$, com corpos fixos $\Q(\sqrt3)$,
$\Q(\sqrt2)$, $\Q(\sqrt6)$ (note que $\sigma\tau$ fixa $\sqrt6 =
\sqrt2\sqrt3$). O reticulado: $\Q$ embaixo, os três corpos
quadráticos no meio, $L$ no topo --- e nada mais
(\cref{thm:b3:galois:fundamental}). Para $(X^2-2)(X^2-3)(X^2-6)$:
o \omterm{thm:b3:galois:splitting}{corpo de decomposição} é o \emph{mesmo} $L$ ($\sqrt6 =
\sqrt2\sqrt3$), de modo que a resposta é idêntica: a
\omterm{thm:b3:galois:fundamental}{correspondência de Galois} é um invariante da \emph{extensão}, não do
polinômio escolhido para apresentá-la.
\end{solution}

\begin{solution}{exo:b3:galois:8}
(a) $G$ age fielmente e transitivamente (irredutibilidade) sobre as
três raízes: $G \hookrightarrow S_3$ com $3 \mid \abs G$: $G
\cong A_3$ ou $S_3$. Todo $\sigma \in G$ permuta os $x_i$, e
$\sigma(\delta) = \varepsilon(\sigma)\,\delta$ ($\delta$ é
alternante nas raízes). Se $\Delta$ é um quadrado em $\Q$:
$\delta \in \Q^\times$ (note $\delta \neq 0$: \omterm{def:b3:galois:separable}{separável}), de modo que
$\varepsilon(\sigma) = 1$ para todo $\sigma$: $G \subseteq A_3$,
logo $= A_3$. Caso contrário: $\delta \notin \Q$, de modo que algum $\sigma$ tem
$\varepsilon(\sigma) = -1$: $G = S_3$. (Em ambos os casos $\Q(\delta)
= L^{G \cap A_3}$.)

(b) Com $x_1 + x_2 + x_3 = 0$: $u + v = 2x_1 - (x_2 + x_3) =
3x_1$. Também
\[
uv = \sum_i x_i^2 + (j + j^2)\sum_{i<k}x_ix_k
= \Bigl(\sum x_i\Bigr)^2 - 3\sum_{i<k}x_ix_k = -3p,
\]
usando $j + j^2 = -1$ e $\sum_{i<k}x_ix_k = p$. Então
\[
u^3 + v^3 = (u+v)^3 - 3uv(u+v) = 27x_1^3 + 9p\cdot 3x_1
= 27\,(x_1^3 + px_1) = -27q .
\]
Assim $u^3, v^3$ são as raízes de $Y^2 + 27qY - 27p^3 = 0$
(produto $(uv)^3 = -27p^3$): $u^3 = \frac{-27q +
\sqrt{729q^2 + 108p^3}}2$, e $x_1 = \frac{u + v}3$ com $v =
-3p/u$: Cardano. A \omterm{def:b3:groups:derived}{solubilidade} de $S_3$ é o esqueleto: a
torre $\Q \subseteq \Q(\delta) \subseteq \Q(\delta, j, u)$
adjunge primeiro uma raiz quadrada ($\delta$, corpo fixo de $A_3$:
o passo $S_3 \to S_3/A_3$), depois uma raiz cúbica ($u$, pois $u^3
\in \Q(\delta, j)$: o passo $A_3 \to \{e\}$) --- a \omterm{def:b3:groups:derived}{série
derivada} $S_3 \supset A_3 \supset \{e\}$ tornada concreta.
\end{solution}

\begin{solution}{exo:b3:galois:9}
$\eta_0 + \eta_1 = \zeta_5 + \zeta_5^2 + \zeta_5^3 + \zeta_5^4 =
-1$ (a soma de todas as raízes $5$-ésimas da unidade é $0$). $\eta_0\eta_1 =
(\zeta + \zeta^4)(\zeta^2 + \zeta^3) = \zeta^3 + \zeta^4 +
\zeta^6 + \zeta^7 = \zeta^3 + \zeta^4 + \zeta + \zeta^2 = -1$
(índices módulo $5$). Logo $\eta_0, \eta_1$ são as raízes de $Y^2 + Y
- 1$: $\frac{-1 \pm \sqrt5}2$. Como $\eta_0 =
2\cos\frac{2\pi}5 > 0$: $\eta_0 = \frac{\sqrt5 - 1}2$, donde
$\cos\frac{2\pi}5 = \frac{\sqrt5 - 1}4$, e $\eta_1 =
\frac{-1-\sqrt5}2$. O grupo
$\operatorname{Gal}(\Q(\zeta_5)/\Q) \cong (\Z/5\Z)^\times$ é
cíclico de ordem $4$: ele tem um \emph{único} subgrupo de ordem
$2$ ($\{\pm 1\}$, isto é, $\zeta \mapsto \zeta^{\pm1}$), de modo que
$\Q(\zeta_5)$ tem um único subcorpo quadrático
(\cref{thm:b3:galois:fundamental}), que contém $\eta_0
\notin \Q$: é $\Q(\eta_0) = \Q(\sqrt5)$. Construtibilidade:
$\cos\frac{2\pi}5$ está na torre quadrática $\Q \subseteq
\Q(\sqrt5)$, e $\zeta_5$ um passo quadrático acima:
o~\cref{thm:b3:galois:wantzel} constrói o pentágono.
\end{solution}

\begin{solution}{exo:b3:galois:10}
(a) Escreva $s = S^{1/p}$, $t = T^{1/p}$ (elementos de um \omterm{def:b3:galois:closure}{fecho
algébrico} escolhido, com $s^p = S$, $t^p = T$). $X^p - S$ é
\omterm{def:b3:rings:divisibility}{irredutível} sobre as frações de $K = \mathbb F_p(S, T) = (\mathbb
F_p(T))(S)$: \omterm{thm:b3:rings:criteria}{Eisenstein} no \omterm{def:b3:rings:divisibility}{elemento primo} $S$ do
\omterm{def:b3:rings:pidufd}{DFU} $\mathbb F_p(T)[S]$ (\cref{thm:b3:rings:criteria}). Logo
$[K(s):K] = p$; do mesmo modo, $X^p - T$ é de \omterm{thm:b3:rings:criteria}{Eisenstein} em $T$ sobre
as frações de $K(s) = \mathbb F_p(s)(T)$ --- $T$ permanece primo em
$\mathbb F_p(s)[T]$ ---, dando $[L : K(s)] = p$ e $[L:K] =
p^2$. Para $\alpha \in L$: $L = K[s, t]$, de modo que $\alpha = \sum
c_{ij}s^it^j$ ($c_{ij} \in K$), e pelo morfismo de \omterm{thm:b3:galois:finitefields}{Frobenius}
$\alpha^p = \sum c_{ij}^p S^iT^j \in K$.

(b) Se $L = K(\alpha)$, então $[K(\alpha):K] = p^2$; mas
$\alpha^p = a \in K$ significa que $\alpha$ anula $X^p - a$, de modo que
$\deg\pi_\alpha \leq p < p^2$: contradição. Nenhum elemento
primitivo: o~\cref{thm:b3:galois:primitive} realmente precisa da
separabilidade (aqui todo $\pi_\alpha$ divide algum $X^p - a =
(X - \alpha)^p$: puramente inseparável).
\end{solution}

\begin{solution}{exo:b3:galois:11}
(a) \omterm{thm:b3:rings:criteria}{Eisenstein} em $2$ ($2 \mid 4, 2$; $4 \nmid 2$): $P$
\omterm{def:b3:rings:divisibility}{irredutível}. Se $\alpha$ é uma raiz, $[\Q(\alpha):\Q] = 5$
divide $[L:\Q] = \abs G$ ($L$ o \omterm{thm:b3:galois:splitting}{corpo de decomposição}): Cauchy
(\cref{thm:b3:groups:cauchy}) dá um elemento de ordem $5$ em
$G \leq S_5$; em $S_5$, só os $5$-ciclos têm ordem $5$ (as ordens
são mmc dos comprimentos dos ciclos).

(b) $P'(x) = 5x^4 - 4$ se anula em $\pm(4/5)^{1/4} \approx \pm
0.946$: um máximo local e depois um mínimo local. Valores: $P(-2)
= -22 < 0$, $P(0) = 2 > 0$, $P(1) = -1 < 0$, $P(2) = 26 > 0$:
três mudanças de sinal, e no máximo três raízes reais (dois pontos
críticos): exatamente $3$ raízes reais, logo um par de raízes
complexas conjugadas. Tome o \omterm{thm:b3:galois:splitting}{corpo de decomposição} $L$ dentro de $\C$:
a conjugação complexa leva $L$ nele mesmo (ela permuta as raízes,
que geram $L$) e fixa $\Q$, de modo que define um elemento de
$G$; ela fixa as três raízes reais e troca as outras duas: uma
transposição.

(c) Sejam $\tau = (a\,b)$ e $\sigma$ um $5$-ciclo em $G$. Alguma
potência $\sigma^k$ leva $a$ em $b$ ($k \ne 0 \bmod 5$), e
$\sigma^k$ é ainda um $5$-ciclo: renomeando, suponha $\sigma =
(1\,2\,3\,4\,5)$ e $\tau = (1\,2)$. Conjugando,
$\sigma^m\tau\sigma^{-m} = (\sigma^m(1)\ \sigma^m(2))$: as
transposições adjacentes $(1\,2), (2\,3), (3\,4), (4\,5),
(5\,1)$ estão todas em $G$; as transposições adjacentes geram $S_5$
(toda transposição $(i\,j)$ é um produto de adjacentes,
e as transposições geram). Logo $G = S_5$, não \omterm{def:b3:groups:derived}{solúvel}
(\cref{cor:b3:groups:snnotsolvable}), e o~\cref{thm:b3:galois:solvable} conclui: $X^5 - 4X + 2$ não é
\omterm{def:b3:galois:radical}{solúvel por radicais}.
\end{solution}

\begin{solution}{exo:b3:galois:12}
(a) $X^4 - 2$ é \omterm{def:b3:rings:divisibility}{irredutível} (\omterm{thm:b3:rings:criteria}{Eisenstein} em $2$):
$[\Q(\alpha):\Q] = 4$; $\iu \notin \Q(\alpha) \subseteq \R$,
de modo que $[L : \Q(\alpha)] = 2$ e $[L:\Q] = 8$. A extensão é
galoisiana (\omterm{thm:b3:galois:splitting}{corpo de decomposição} de um \omterm{def:b3:galois:separable}{polinômio separável}: as raízes
são $\iu^k\alpha$), logo $\abs G = 8$. Um automorfismo leva
$\alpha$ a uma das quatro raízes e $\iu$ a $\pm\iu$:
no máximo $8$ aplicações, todas realizadas. O $\sigma$ indicado (de ordem
$4$: $\sigma^2(\alpha) = -\alpha$, $\sigma^4 = e$) e $\tau$
(de ordem $2$) satisfazem
\[
\tau\sigma\tau(\alpha) = \tau\sigma(\alpha) =
\tau(\iu\alpha) = -\iu\alpha = \sigma^{-1}(\alpha),
\qquad \tau\sigma\tau(\iu) = \iu\cdot(-1)(-1) = \iu ,
\]
mais precisamente: $\tau\sigma\tau(\iu) = \tau\sigma(-\iu) =
\tau(-\iu) = \iu = \sigma^{-1}(\iu)$. Logo $\tau\sigma\tau =
\sigma^{-1}$: a apresentação de $D_4$.

(b) Os dez subgrupos de $D_4 = \langle\sigma,
\tau\rangle$: $\{e\}$; cinco de ordem $2$:
$\langle\sigma^2\rangle$, $\langle\tau\rangle$,
$\langle\sigma^2\tau\rangle$, $\langle\sigma\tau\rangle$,
$\langle\sigma^3\tau\rangle$; três de ordem $4$:
$\langle\sigma\rangle$, $\{e, \sigma^2, \tau,
\sigma^2\tau\}$, $\{e, \sigma^2, \sigma\tau,
\sigma^3\tau\}$; e $D_4$. Corpos fixos (\omterm{def:b3:galois:extension}{grau} = índice):
$\{e\} \leftrightarrow L$; os subgrupos de ordem $2$
$\leftrightarrow$ os cinco corpos quárticos
\[
\langle\tau\rangle \leftrightarrow \Q(\alpha),\quad
\langle\sigma^2\tau\rangle \leftrightarrow \Q(\iu\alpha),
\quad
\langle\sigma^2\rangle \leftrightarrow \Q(\sqrt2, \iu),\quad
\langle\sigma\tau\rangle \leftrightarrow
\Q\bigl((1+\iu)\alpha\bigr),\quad
\langle\sigma^3\tau\rangle \leftrightarrow
\Q\bigl((1-\iu)\alpha\bigr) .
\]
Verificações: $\tau$ fixa o real $\alpha$; $\sigma^2\tau$ leva
$\alpha \mapsto -\alpha$ e $\iu \mapsto -\iu$, fixando
$\iu\alpha$; e como $\sigma\tau(\alpha) = \iu\alpha$,
$\sigma\tau(\iu) = -\iu$:
\[
\sigma\tau\bigl((1+\iu)\alpha\bigr) =
(1 - \iu)\,\iu\alpha = (1 + \iu)\alpha,
\qquad
\sigma^3\tau\bigl((1-\iu)\alpha\bigr) =
(1 + \iu)(-\iu)\alpha = (1 - \iu)\alpha :
\]
cada reflexão fixa seu gerador, e o corpo fixo,
de \omterm{def:b3:galois:extension}{grau} $4 = $ índice, é exatamente o corpo que ele gera
(o gerador é raiz de $X^4 + 8$, \omterm{def:b3:rings:divisibility}{irredutível}). Os subgrupos de ordem $4$
$\leftrightarrow$ os três corpos quadráticos:
$\langle\sigma\rangle \leftrightarrow \Q(\iu)$ ($\sigma$
fixa $\iu$); $\{e, \sigma^2, \tau, \sigma^2\tau\}
\leftrightarrow \Q(\sqrt2)$ (todos os quatro fixam $\alpha^2$ a menos de
verificações de sinal: $\tau(\sqrt2) = \sqrt2$, $\sigma^2(\alpha^2) =
(-\alpha)^2$); $\{e, \sigma^2, \sigma\tau, \sigma^3\tau\}
\leftrightarrow \Q(\iu\sqrt2)$ ($\sigma\tau(\iu\alpha^2) =
(-\iu)(\iu\alpha)^2 = \iu\alpha^2$).

(c) Galoisiano sobre $\Q$ $\leftrightarrow$ \omterm{def:b3:groups:normal}{subgrupos normais} de
$D_4$: $\{e\}$, $\langle\sigma^2\rangle$ (o \omterm{ex:b3:groups:actions}{centro}), os
três subgrupos de ordem $4$ e $D_4$ --- de modo que os corpos
intermediários galoisianos são $L$, $\Q(\sqrt2, \iu)$, os três
corpos quadráticos e $\Q$. Os cinco corpos quárticos fixados por
reflexões não normais não são galoisianos: $\Q(\alpha)$ contém
uma raiz de $X^4 - 2$, mas não $\iu\alpha$ (ele é real) ---
a conjugação por $\sigma$ move $\langle\tau\rangle$ para
$\langle\sigma^2\tau\rangle$, exatamente como move
$\Q(\alpha)$ para $\Q(\iu\alpha)$: a não normalidade do
subgrupo \emph{é} a existência de um corpo conjugado.
\end{solution}

\begin{solution}{pb:b3:galois:1}
\textbf{1.} $\Phi_{17}$ é \omterm{def:b3:rings:divisibility}{irredutível}
(\cref{thm:b3:galois:cyclotomicirred}, ou o~\cref{ex:b3:rings:eisenstein} para índice primo): $[L:\Q] =
\varphi(17) = 16$ e $G \cong (\Z/17\Z)^\times$, cíclico de
ordem $16$ (\cref{thm:b3:galois:cyclic}). Potências de $3$ módulo
$17$:
\[
3,\ 9,\ 10,\ 13,\ 5,\ 15,\ 11,\ 16,\ 14,\ 8,\ 7,\ 4,\ 12,\ 2,\
6,\ 1
\]
--- dezesseis valores distintos: $3$ gera.

\textbf{2.} $G = \langle\sigma\rangle$ cíclico de ordem $16$;
$H_k = \langle\sigma^{2^k}\rangle$ tem ordem $2^{4-k}$, e
$[H_k : H_{k+1}] = 2$. Pelo teorema fundamental
(\cref{thm:b3:galois:fundamental}), $L_k = L^{H_k}$ satisfazem
$[L_k : \Q] = [G : H_k] = 2^k$: cada $[L_{k+1}:L_k] = 2$.

\textbf{3.} $\zeta \in L = L_4$ está no topo de uma torre de extensões
quadráticas de $\Q$: pelo~\cref{thm:b3:galois:wantzel}, $\zeta$ é
\omterm{def:b3:galois:constructible}{construtível}; o $17$-ágono tem vértices $\zeta^k$.

\textbf{4.} $\sigma^2$ multiplica os expoentes por $9$; os
expoentes de $\eta_0$ são as potências pares de $3$,
\[
\{3^{2k} \bmod 17\} = \{1, 9, 13, 15, 16, 8, 4, 2\},
\]
um conjunto estável pela multiplicação por $9 = 3^2$; assim $\eta_0$
(e, do mesmo modo, $\eta_1$) é fixado por $H_1 =
\langle\sigma^2\rangle$: $\eta_0, \eta_1 \in L_1$, um corpo
quadrático. $\sigma$ leva potências pares em ímpares: ele troca $\eta_0,
\eta_1$. Logo $\eta_0 + \eta_1$ e $\eta_0\eta_1$ são fixados por
todo $G$: racionais; $\eta_0, \eta_1$ são as raízes de uma
quadrática racional.

\textbf{5.} $\eta_0 + \eta_1 = \sum_{c=1}^{16}\zeta^c = -1$. O
produto se expande em $64$ termos $\zeta^{a + b}$, com $a$ no conjunto par
e $b$ no conjunto ímpar. Nenhum termo é $\zeta^0$: $b = -a$ é
impossível, pois $-1 = 16 = 3^8$ é uma potência \emph{par}, de modo que
$-a$ permanece no conjunto par. Assim $\eta_0\eta_1 = \sum_{c \neq 0}
n_c\zeta^c$ com $\sum n_c = 64$; aplicar $\sigma$ fixa
$\eta_0\eta_1$ (ele troca os fatores) e permuta os $\zeta^c$
transitivamente sobre todos os $c \neq 0$, de modo que todos os $n_c$ são iguais: $n_c
= 4$ e $\eta_0\eta_1 = 4\sum_{c\neq0}\zeta^c = -4$.

\textbf{6.} $\eta_{0,1}$ resolvem $Y^2 + Y - 4 = 0$: $\frac{-1 \pm
\sqrt{17}}2$. Numericamente, agrupando expoentes conjugados,
$\eta_0 = 2\bigl(\cos\tfrac{2\pi}{17} + \cos\tfrac{4\pi}{17} +
\cos\tfrac{8\pi}{17} + \cos\tfrac{16\pi}{17}\bigr) \approx 1.56
> 0$: $\eta_0 = \frac{-1+\sqrt{17}}2$, $\eta_1 =
\frac{-1-\sqrt{17}}2$, e $L_1 = \Q(\eta_0) = \Q(\sqrt{17})$.

\textbf{7.} Os conjuntos de expoentes: $\beta_0$: $\{1, 13, 16, 4\}$ =
potências $3^{4k}$; $\beta_2$: $\{9, 15, 8, 2\}$ = $9 \times$ desse
conjunto. União: o conjunto par: $\beta_0 + \beta_2 = \eta_0$; do mesmo modo,
$\beta_1 + \beta_3 = \eta_1$. A multiplicação por $13 = 3^4$
estabiliza o conjunto de expoentes de cada $\beta_i$: fixado por $H_2 =
\langle\sigma^4\rangle$; e $\sigma^2$ ($\times 9$) leva
$\{1,13,16,4\}$ em $\{9, 15, 8, 2\}$: troca $\beta_0,
\beta_2$.

\textbf{8.} Expandindo $\beta_0\beta_2$, as dezesseis somas de
expoentes
\[
\{1,13,16,4\} + \{9,15,8,2\} =
\{10,16,9,3,\ 5,11,4,15,\ 8,14,7,1,\ 13,2,12,6\}
\]
cobrem $1, \dots, 16$ exatamente uma vez: $\beta_0\beta_2 =
\sum_{c\ne0}\zeta^c = -1$. Aplicando $\sigma$ (que leva
$\beta_0 \mapsto \beta_1$, $\beta_2 \mapsto \beta_3$:
expoentes $\times 3$): $\beta_1\beta_3 = \sigma(\beta_0\beta_2)
= -1$.

\textbf{9.} $\beta_0, \beta_2$ resolvem $Y^2 - \eta_0 Y - 1 = 0$,
de modo que $\beta_0 = \frac{\eta_0 + \sqrt{\eta_0^2 + 4}}2$ (numericamente
$\beta_0 = 2\cos\frac{2\pi}{17} + 2\cos\frac{8\pi}{17} \approx
2.05 > 0$, o sinal $+$). Do mesmo modo, $\beta_1 = \frac{\eta_1 +
\sqrt{\eta_1^2 + 4}}2 \approx 0.344$ (a verificação numérica fixa o
sinal de novo). $L_2 = L_1(\beta_0)$, quadrático sobre $L_1$.

\textbf{10.} $\gamma_0 + \gamma_1 = \zeta + \zeta^{16} +
\zeta^{13} + \zeta^4 = \beta_0$. E
\[
\gamma_0\gamma_1 = (\zeta + \zeta^{16})(\zeta^{13} + \zeta^4)
= \zeta^{14} + \zeta^{5} + \zeta^{12} + \zeta^{3} = \beta_1 .
\]
Logo $\gamma_0, \gamma_1$ resolvem $Y^2 - \beta_0Y + \beta_1 = 0$;
numericamente $\gamma_0 = 2\cos\frac{2\pi}{17} \approx 1.865 >
\gamma_1 \approx 0.185$: $\gamma_0 = \frac{\beta_0 +
\sqrt{\beta_0^2 - 4\beta_1}}2$.

\textbf{11.} Encadeando:
\[
\eta_0 = \frac{-1 + \sqrt{17}}2,
\quad
\beta_0 = \frac{\eta_0 + \sqrt{\eta_0^2 + 4}}2,
\quad
\beta_1 = \frac{\eta_1 + \sqrt{\eta_1^2 + 4}}2,
\quad
\cos\frac{2\pi}{17} = \frac{\beta_0 + \sqrt{\beta_0^2 -
4\beta_1}}4 .
\]
Numericamente: $\sqrt{17} \approx 4.1231$, $\eta_0 \approx
1.5616$, $\eta_1 \approx -2.5616$, $\beta_0 \approx 2.0494$,
$\beta_1 \approx 0.3441$, $\beta_0^2 - 4\beta_1 \approx
2.8234$ e $\cos\frac{2\pi}{17} \approx \frac{2.0494 +
1.6803}4 \approx 0.93242$ --- contra $\cos\frac{2\pi}{17} =
0.93247\dots$: a pequena discrepância é arredondamento nas
fórmulas intermediárias; carregar mais dígitos reproduz
$0.932472$.

\textbf{12.} A construção exigiu que $[L:\Q] = p - 1$ fosse uma
potência de $2$, para que exista uma cadeia completa de subgrupos
de índice $2$. Se $p = 2^m + 1$ é primo e $m = ab$ com $a$ ímpar
$>
1$: $x + 1 \mid x^a + 1$ em $x = 2^b$ mostra que $2^b + 1$ divide
propriamente $p$ --- impossível. Logo $m$ é uma potência de
$2$: $p =
2^{2^t} + 1$, um \emph{primo de Fermat} ($3, 5, 17, 257, 65537$,
\dots). Reciprocamente, para tais $p$, $\varphi(p) = 2^{2^t}$ e o
argumento das questões 1--3 (ou
o~\cref{cor:b3:galois:impossible}) se aplica: o $p$-ágono regular
é \omterm{def:b3:galois:constructible}{construtível} se, e somente se, $p$ é um primo de Fermat.

\textbf{13.} $\varphi(n)$ é uma potência de $2$ exatamente quando $n =
2^a p_1\cdots p_r$ com primos de Fermat $p_i$ distintos
(multiplicatividade de $\varphi$; uma potência de primo ímpar $p^k$, $k
\geq 2$, contribui com o fator $p \nmid 2^m$). Para $n \leq
20$, os $n$-ágonos regulares \omterm{def:b3:galois:constructible}{construtíveis} são
\[
n = 3, 4, 5, 6, 8, 10, 12, 15, 16, 17, 20
\]
com os valores respectivos
\[
\varphi(n) = 2,\ 2,\ 4,\ 2,\ 4,\ 4,\ 4,\ 8,\ 8,\ 16,\ 8 .
\]
Os impossíveis são $n = 7, 9, 11, 13, 14, 18, 19$, para os quais
$\varphi(n) = 6, 6, 10, 12, 6, 6, 18$ tem um fator primo ímpar.

\textbf{14.} Os quadrados formam a imagem do morfismo de elevação ao
quadrado no cíclico $(\Z/p\Z)^\times$, de índice $2$:
$\frac{p-1}2$ quadrados, $\frac{p-1}2$ não quadrados, de modo que os
símbolos somam $0$. Então
\[
\sum_{a=0}^{p-1}\zeta^{a^2} = 1 + 2\sum_{b \text{ quadrado}
\neq 0}\zeta^b
= 1 + \sum_{b\neq0}\Bigl(1 + \Bigl(\frac bp\Bigr)\Bigr)
\zeta^b = \sum_{b}\zeta^b + g = g,
\]
usando $1 + \bigl(\frac bp\bigr) = \#\{a : a^2 = b\}$ e
$\sum_{b=0}^{p-1}\zeta^b = 0$. Para $p = 17$: os quadrados módulo
$17$ são as potências pares do gerador $3$, isto é, os
expoentes que aparecem em $\eta_0$ (Parte II), de modo que $g =
\sum_{\text{par }k}\zeta^{3^k} - \sum_{\text{ímpar }k}
\zeta^{3^k} = \eta_0 - \eta_1$.

\textbf{15.} Com $b = c - a$ ($a, b$ percorrem os resíduos
não nulos, $c = a + b$ todos os resíduos):
\[
g^2 = \sum_c\zeta^c\sum_{a\neq0,\,a\neq c}
\Bigl(\frac{a(c-a)}p\Bigr) .
\]
Para $c = 0$: $\bigl(\frac{-a^2}p\bigr) =
\bigl(\frac{-1}p\bigr)$, somado sobre $p - 1$ valores. Para $c \neq 0$: substitua $c - a = at$, isto é, $t = c/a -
1$; quando $a$ percorre os resíduos não nulos, $t$ percorre
bijetivamente os resíduos $\neq -1$ (inverta: $a = c/(1 +
t)$). O somando se torna $\bigl(\frac{a^2t}p\bigr) =
\bigl(\frac tp\bigr)$, e
\[
\sum_{t \neq -1}\Bigl(\frac tp\Bigr)
= -\Bigl(\frac{-1}p\Bigr)
\]
(a soma completa se anula pela questão 14). Logo
\[
g^2 = \Bigl(\frac{-1}p\Bigr)\Bigl[(p-1) -
\sum_{c\neq0}\zeta^c\Bigr]
= \Bigl(\frac{-1}p\Bigr)\,p = p^* ,
\]
usando $\sum_{c\neq0}\zeta^c = -1$ e o critério de Euler
$\bigl(\frac{-1}p\bigr) = (-1)^{(p-1)/2}$. Para $p = 17$:
$(\eta_0 - \eta_1)^2 = (\eta_0 + \eta_1)^2 - 4\eta_0\eta_1 =
1 + 16 = 17$, de acordo com a Parte II.

\textbf{16.} $g^2 = p^*$ exibe $\sqrt{p^*} = \pm g \in
\Q(\zeta_p)$, de modo que $\Q(\sqrt{p^*})$ é um subcorpo quadrático.
Unicidade: os subcorpos de \omterm{def:b3:galois:extension}{grau} $2$ correspondem, pela
\omterm{thm:b3:galois:fundamental}{correspondência de Galois}, aos subgrupos de índice $2$ do cíclico
$\operatorname{Gal}(\Q(\zeta_p)/\Q) \cong (\Z/p\Z)^\times$,
e um grupo cíclico de ordem par tem exatamente um tal
subgrupo (os quadrados). Todo corpo quadrático é
$\Q(\sqrt{d})$ com $d$ livre de quadrados, e combinar os corpos
$\Q(\sqrt{p^*})$, $\Q(\iu) \subseteq \Q(\zeta_4)$ e
$\Q(\sqrt2) \subseteq \Q(\zeta_8)$ dentro de um mesmo
$\Q(\zeta_N)$ alcança todo $\sqrt d$: o caso quadrático de
Kronecker--Weber.

\textbf{17.} Em qualquer anel comutativo, $(x + y)^q = x^q + y^q
+ q(\cdots)$: os coeficientes binomiais $\binom qk$, $0 < k <
q$, são divisíveis pelo primo $q$. Iterando sobre os $p - 1$
termos de $g$:
\[
g^q \equiv \sum_{a}\Bigl(\frac ap\Bigr)^{q}\zeta^{aq}
= \sum_a\Bigl(\frac ap\Bigr)\zeta^{aq}
\pmod{q\Z[\zeta]},
\]
($q$ ímpar: o símbolo não muda). Reindexe $b = aq$: $a =
q^{-1}b$ e $\bigl(\frac{q^{-1}b}p\bigr) =
\bigl(\frac{q}p\bigr)\bigl(\frac bp\bigr)$ (multiplicatividade;
$\bigl(\frac{q^{-1}}p\bigr) = \bigl(\frac qp\bigr)$, pois o
símbolo de um inverso é igual ao símbolo): $g^q \equiv
\bigl(\frac qp\bigr)g$.

\textbf{18.} $g^q = g\,(g^2)^{(q-1)/2} = g\,(p^*)^{(q-1)/2}$
exatamente (questão 15), e o critério de Euler em $\Z$ dá
$(p^*)^{(q-1)/2} \equiv \bigl(\frac{p^*}q\bigr) \pmod q$,
logo módulo $q\Z[\zeta]$: $g^q \equiv \bigl(\frac{p^*}q\bigr)g$.
Comparando com a questão 17 e multiplicando por $g$:
\[
\Bigl(\frac qp\Bigr)p^* \equiv \Bigl(\frac{p^*}q\Bigr)p^*
\pmod{q\Z[\zeta]} .
\]
Ambos os lados são inteiros racionais; sua diferença, $0$ ou
$\pm2p^*$, está em $q\Z[\zeta] \cap \Z = q\Z$ (um inteiro $m
\in q\Z[\zeta]$ tem $m/q \in \Q \cap \Z[\zeta] = \Z$, esta última porque $1, \zeta, \dots, \zeta^{p-2}$ é uma
$\Q$-base com coordenadas racionais das quais se lê a
integralidade). Como $q \nmid 2p^*$ ($q$ ímpar, $q \neq p$),
a diferença é $0$: $\bigl(\frac qp\bigr) =
\bigl(\frac{p^*}q\bigr)$.

\textbf{19.} Por multiplicatividade, $\bigl(\frac{p^*}q\bigr) =
\bigl(\frac{-1}q\bigr)^{(p-1)/2}\bigl(\frac pq\bigr) =
(-1)^{\frac{q-1}2\cdot\frac{p-1}2}\bigl(\frac pq\bigr)$,
de modo que a questão 18 se lê $\bigl(\frac qp\bigr)\bigl(\frac
pq\bigr) = (-1)^{\frac{p-1}2\frac{q-1}2}$: a reciprocidade.
Verificação $(17, 3)$: o expoente $\frac{16}2\cdot\frac22 = 8$
é par, de modo que os dois símbolos devem coincidir; os quadrados módulo $3$ são $\{1\}$ e $17 \equiv 2$:
$\bigl(\frac{17}3\bigr) = -1$; os quadrados módulo $17$ são $\{1, 4,
9, 16, 8, 2, 15, 13\}$ e $3$ está ausente:
$\bigl(\frac3{17}\bigr) = -1$. Produto $+1$, como previsto. Para $x^2
\equiv 219 \pmod{383}$: $\bigl(\frac{219}{383}\bigr) =
\bigl(\frac3{383}\bigr)\bigl(\frac{73}{383}\bigr)$. Primeiro:
$383 \equiv 3 \pmod4$ e $3 \equiv 3$: a reciprocidade dá
$\bigl(\frac3{383}\bigr) = -\bigl(\frac{383}3\bigr) =
-\bigl(\frac23\bigr) = -(-1) = +1$. Segundo: $73 \equiv 1
\pmod 4$: $\bigl(\frac{73}{383}\bigr) =
\bigl(\frac{383}{73}\bigr) = \bigl(\frac{18}{73}\bigr) =
\bigl(\frac2{73}\bigr)$ ($18 = 2\cdot3^2$), e $73 \equiv 1
\pmod 8$ torna $2$ um quadrado módulo $73$ (lei complementar,
demonstrável por $g = \zeta_8 + \zeta_8^{-1} = \sqrt2$ em
$\Q(\zeta_8)$ pelo mesmo método): $+1$. Total $+1$: a
congruência é \omterm{def:b3:groups:derived}{solúvel}.

\textbf{20.} Existência e unicidade da raiz primitiva:
se $w$ tem conjunto de períodos $\{d : w = u^{n/d},\ \abs u = d\}$,
o menor tal $d_0$ divide todo outro período $d$
(se $w$ é ao mesmo tempo uma $d$-potência e uma $d'$-potência, é uma
$\gcd(d, d')$-potência: compare letras em índices congruentes
módulo o mdc, via Bézout), e o bloco de comprimento $d_0$ é
primitivo. Ordenando as $q^n$ palavras pelo comprimento de sua
raiz primitiva: $q^n = \sum_{d \mid n}A(d)$. Como $A$ e
$d\,N_q(d)$ satisfazem a mesma recorrência com os mesmos valores
para $n = 1$ (cada uma determina a outra indutivamente a partir de $q^n
= \sum_{d\mid n}(\cdot)$), elas são iguais: $A(d) =
d\,N_q(d)$. Bijeção: um elemento $\alpha$ de \omterm{def:b3:galois:extension}{grau} $d$
fornece a palavra $w_\alpha$ dos coeficientes de\dots\
melhor, diretamente: os elementos de \omterm{def:b3:galois:extension}{grau} $d$ em
$\overline{\mathbb F_q}$ são as raízes dos $N_q(d)$
irredutíveis de \omterm{def:b3:galois:extension}{grau} $d$, cada um contribuindo com suas $d$
raízes distintas (separabilidade): $d\,N_q(d)$ elementos de
\omterm{def:b3:galois:extension}{grau} $d$, coincidindo com a contagem $q^n = \sum_{d\mid n}\#\{
\text{elementos de grau } d \text{ em } \mathbb F_{q^n}\}$
--- o mesmo crivo, uma vez sobre palavras, outra sobre elementos do
corpo.

\textbf{21.} Lema: $\sum_{d \mid m}\mu(d) = \mathbf
1_{m=1}$. Para $m > 1$, fixe um primo $p \mid m$: os divisores livres de quadrados
de $m$ se emparelham como $\{d, pd\}$ com $p \nmid d$,
e $\mu(pd) = -\mu(d)$: a soma se cancela. Então, para $f(n) =
\sum_{e\mid n}g(e)$:
\[
\sum_{d \mid n}\mu(d)\,f\bigl(\tfrac nd\bigr)
= \sum_{d \mid n}\mu(d)\!\!\sum_{e \mid n/d}\!\!g(e)
= \sum_{e \mid n}g(e)\!\!\sum_{d \mid n/e}\!\!\mu(d)
= g(n) .
\]
Com $f(n) = q^n$ e $g(n) = nN_q(n)$
(\cref{exo:b3:galois:6}): $N_q(n) =
\frac1n\sum_{d\mid n}\mu(d)\,q^{n/d}$.

\textbf{22.} O termo $d = 1$ é $q^n$; todo outro termo tem
$\abs{\mu(d)q^{n/d}} \leq q^{n/2}$, e grosseiramente
$\sum_{d \mid n, d > 1}q^{n/d} \leq \sum_{j \leq
n/2}q^j < 2q^{n/2}$ (geométrica, $q \geq 2$). Logo $nN_q(n) >
q^n - 2q^{n/2} \geq 0$ para $n \geq 1$: existem irredutíveis de todo
\omterm{def:b3:galois:extension}{grau}, e $\mathbb F_q[X]/(\pi) = \mathbb F_{q^n}$ é
(re)construído --- existência com um recenseamento. A proporção de
irredutíveis entre os $q^n$ polinômios mônicos de \omterm{def:b3:galois:extension}{grau} $n$ é
$\frac1n(1 + O(q^{-n/2}))$: o teorema dos números primos de
$\mathbb F_q[X]$, com $n$ no papel de $\log x$. Para $q = 2$, as
contagens $2, 1, 2, 3$ do~\cref{exo:b3:galois:6} batem com a
fórmula: por exemplo, $N_2(4) = \frac14(2^4 - 2^2) = 3$.

\textbf{23.} No grupo abeliano multiplicativo das funções racionais
não nulas sobre $\mathbb F_q$, ponha $F(n) = X^{q^n} -
X$ e $G(n) = \prod_{\deg\pi = n}\pi$;
o~\cref{exo:b3:galois:6} diz que $F(n) = \prod_{d\mid n}G(d)$.
O argumento de Möbius da questão 21, escrito
multiplicativamente (os expoentes se somam exatamente como as somas
faziam), dá $G(n) = \prod_{d \mid n}F(d)^{\mu(n/d)}$. Para $q = 2$,
$n = 2$: $G(2) = \frac{X^4 - X}{X^2 - X} = \frac{X(X^3 -
1)}{X(X - 1)} = X^2 + X + 1$, a única quadrática \omterm{def:b3:rings:divisibility}{irredutível}
sobre $\mathbb F_2$, como deve ser.

\textbf{24.} $\omega^2 = \iu$ e $\omega^{-2} = -\iu$, de modo que
$g^2 = \omega^2 + 2 + \omega^{-2} = 2$. Sonho do calouro no
anel comutativo $\Z[\omega]/q\Z[\omega]$: $g^q =
(\omega + \omega^{-1})^q \equiv \omega^q + \omega^{-q}$. O
valor de $\omega^q + \omega^{-q}$ depende apenas de $q \bmod 8$:
para $q \equiv \pm1$, $\omega^q + \omega^{-q} = \omega^{\pm1} +
\omega^{\mp1} = g$; para $q \equiv \pm3$, usando $\omega^4 =
-1$, $\omega^{3} = -\omega^{-1}$ e $\omega^{-3} = -\omega$,
de modo que $\omega^q + \omega^{-q} = -g$. Por outro lado, $g^q =
g\,(g^2)^{(q-1)/2} = g\,2^{(q-1)/2} \equiv \bigl(\frac2q\bigr)
g \pmod{q\Z[\omega]}$ pelo critério de Euler módulo $q$. Comparando
e multiplicando por $g$: $2\bigl(\frac2q\bigr) \equiv
\pm2 \pmod{q\Z[\omega]}$; se os sinais discordassem, $q$
dividiria $4$ em $\Z[\omega]$, logo em $\Z$ ($q\Z[\omega] \cap
\Z = q\Z$: coordenadas na base $1, \omega, \omega^2,
\omega^3$), impossível para $q$ ímpar. Logo $\bigl(\frac2q\bigr) =
+1$ se, e somente se, $q \equiv \pm1 \pmod 8$. Verificação de paridade: $q = 8k \pm 1$
dá $(q^2 - 1)/8 = 2k(4k \pm 1)$, par; $q = 8k \pm 3$
dá $(q^2 - 1)/8 = 8k^2 \pm 6k + 1$, ímpar: a fórmula
$(-1)^{(q^2-1)/8}$ codifica a separação em casos. Numericamente: $3^2
= 9 \equiv 2 \pmod 7$ ($7 \equiv -1$), $6^2 = 36 \equiv 2
\pmod{17}$ ($17 \equiv 1$); os quadrados módulo $3$ são $\{0,
1\}$ e módulo $5$ são $\{0, 1, 4\}$, e nenhum contém $2$
($3 \equiv 3$, $5 \equiv -3 \pmod 8$).

\textbf{25.} Todo $f \in \mathbb F_q[X]$ mônico se fatora
unicamente como $\prod_\pi\pi^{e_\pi}$ sobre os mônicos
irredutíveis: ordenando por \omterm{def:b3:galois:extension}{grau},
\[
\sum_{f \text{ mônico}}t^{\deg f}
= \prod_\pi\ \sum_{e \geq 0}t^{e\deg\pi}
= \prod_\pi\bigl(1 - t^{\deg\pi}\bigr)^{-1}
= \prod_{n \geq 1}\bigl(1 - t^n\bigr)^{-N_q(n)},
\]
todos os produtos legítimos t-adicamente (apenas os \omterm{def:b3:galois:extension}{graus} $\leq m$
afetam o coeficiente de $t^m$, e há finitos
irredutíveis de cada \omterm{def:b3:galois:extension}{grau}). O lado esquerdo é $\sum_m
q^mt^m = (1 - qt)^{-1}$: a identidade. Logaritmos: $-\log(1 -
qt) = \sum_m\frac{q^m}mt^m$, ao passo que $\sum_nN_q(n)\bigl(-\log(1
- t^n)\bigr) = \sum_nN_q(n)\sum_k\frac{t^{nk}}k$; o
coeficiente de $t^m$ dá $\frac{q^m}m = \sum_{dk = m}
\frac{N_q(d)}k = \frac1m\sum_{d \mid m}d\,N_q(d)$, isto é,
$q^m = \sum_{d\mid m}d\,N_q(d)$. Verificação à mão, $q = 2$,
coeficiente de $t^2$: $N_2(1) = 2$, $N_2(2) = 1$, e
$(1 - t)^{-2}(1 - t^2)^{-1} = (1 + 2t + 3t^2 + \dots)(1 + t^2
+ \dots)$ tem coeficiente de $t^2$ igual a $3 + 1 = 4 = 2^2$. Finalmente,
a fórmula da questão 21 com os divisores $1, 2, 3, 6$:
\[
N_2(6) = \tfrac16\bigl(2^6 - 2^3 - 2^2 + 2\bigr)
= \tfrac{54}6 = 9,
\]
e, de fato, $1\cdot2 + 2\cdot1 + 3\cdot2 + 6\cdot9 = 2 + 2 +
6 + 54 = 64 = 2^6$.
\end{solution}
